% ------------------------------------------------------------
% Page 99
% ------------------------------------------------------------

Je dormis mal. « le temple », « La sacristie », et « c’était trop dur pour elle » me
donnèrent des cauchemars ; à quoi devais-je m’attendre ?

Je fus matinal, et c’est accompagné d’Adrien ROUSSET que j’entre par la sacristie ;
une visite rapide de la salle de classe, que je juge petite pour un Temple mais bien
vaste pour quelques élèves. Quatre longues tables, à bancs mobiles, témoignent d’une
fréquentation qui fut importante. De toute façon je m’en accommoderai ; l’état de la
sacristie, mon futur logement me préoccupe davantage..

La cuisine, toute en longueur sur la largeur du temple, quoique réduite par la montée
de l’escalier, ma paraît aménageable. La chambre ! du haut de l’escalier, je découvre
le désastre. La fenêtre, qui a dû battre tout l’été, a perdu toutes ses vitres, dont les
débris jonchent toute la pièce ; et le cadre lui-même semble difficilement rafistolable.
Il n’y a pas de volets : il n’y en a jamais eu !
-« Quelle histoire !, dit Adrien, Cette fenêtre doit être remplacée. Il faut voir le
Maire, qui ne veut rien entendre, dès qu’il s’agit de l’école ».

La chambre, elle-même est toute petite, fortement rétrécie par une sorte de garde-fou
qui, prudemment, la sépare de l’escalier.
Je suis de plus en plus contrarié, et, au cours de la présentation de ma nomination à
monsieur le Maire, je demanderai le remplacement de la fenêtre.
Il n’y a pas de cour. La cour c’est la ruelle. Mais où sont donc les WC ? Il n’y en a
pas et il n’y en a jamais eu !!

Je ne fais pas de commentaire ; Je les réserve à Monsieur le Maire.

J’erre dans le hameau. La dizaine de familles qui y vivent encore atteste de son
importance. Mais peu de jeunes. Je me présente au hasard des rencontres. Et me voici
chez Mme Argeliès, une accorte personne, très accueillante et fort bavarde. C’est
avec plaisir qu’elle accepte de me prendre en pension. Je lui dis qu’après déjeuner je
dois descendre à Meyrueis me présenter au Maire.
« Ah, je vois ! Comme tous ceux qui vous ont précédé, vous allez discuter avec le
maire comme la pauvre demoiselle de l’an passé. Elle en est revenue malade… Le
soir elle n’a rien pu avaler ;… »
et Ricou, son fils, qui me propose spontanément son vélo, d’ajouter :
« Les maires de Meyrueis n’aiment pas les Oubrets, et celui-là, particulièrement,
parce que nous sommes tous protestants, et lui, c’est « un curotin » ! nous ne votons
pas pour lui ! »

J’apprécie cette liberté de langage, mais j’avais déjà tout compris, effaré de constater
que puissent encore s’exprimer, dans ce lieu, des passions d’un autre temps.

Je passe le reste de la matinée à préparer, à ruminer, à prévoir tous les déroulements
possibles de l’entretien. Le mot « temple » ne sortira pas de ma bouche car je suis
instituteur laïque, bien décidé à rappeler à Monsieur le Maire ses devoirs, et les
responsabilités qui, légalement, sont les siennes.
